\cotefile{158}
\scanpage{71}\mspage{\textcircled{6}}
\diagram{\begin{tikzcd}[ampersand replacement=\&,row sep=0pt]
 F \& \& \\
 X_{\deleted{\unsure{/F|/f}}}\arrow[r,hook,"i"'] \& X' \& e\arrow[l,hook]
\end{tikzcd}}
\[
 X_{/e}\simeq X_{/F}\qquad F\ \unsure{0|o}\text{-asphérique}\ \unsure{\Leftarrow|\Longleftrightarrow}\ i\ \text{cocofinal}
\]
Si $X^{\unsure{\circ|0|o}}$ \unsure{t.t.|tot.|est} $W$-asph., donc \inserted{$X'^{\circ}$ donc} $X'$ est $W$ asph., donc $i\in W$, et si $W$ satisfait $\mathrm{Loc}(6)$, $i$ est $W$-asph., i.e. $X^{\unsure{\prime}}_{/e}$ $W$-asph. i.e. $X_{/F}$ $W$ asph.

Inversement, supposons que tt objet $0$-connexe $F$ de $X^\wedge$ soit $W$-asph. Soit
\[
 X\longrightarrow Y\qquad\text{cocofinal, i.e. les }X_{/y}\ 0\text{-conn.}
\]
\deleted{Il s'ensuit que pour \unsure{tt}} \inserted{Mais $X_{/y}$ est étale sur $X$,} \deleted{\unsure{$Y'$} étale sur $Y$} \inserted{et \unsure{étant|ét. x} $0$-conn.,} il est $W$-asph. Donc $X\to Y$ est $W$-asph. donc $\in W$, donc $Y$ est $W$-asph. Ainsi :

\uline{Prop} \deleted{Si $W$ satisfait \unsure{L}} Pour que $X^\circ$ soit \unsure{tot.|t.t.} \deleted{\unsure{as}} $W$-asph., il \deleted{\illegible{1}} s. que tt objet $0$-connexe de $X^\wedge$ soit $W$-asph., et elle est aussi nécessaire si $W$ \deleted{\unsure{Loc}} satisfait $\mathrm{Loc}(6)$,

\uline{Cor.} (Pour mémoire) \deleted{\unsure{on}} Si \deleted{$X^\circ$ est \unsure{cofiltrante|cofiltrant} et} $W$ satisfait $\mathrm{Loc}(8)$, alors pour \unsure{tte} $X$ telle que $X^\circ$ soit filtrante (i.e. $X$ cofiltrante), les objets $0$-connexes \inserted{$F$} de $X^\wedge$ \deleted{soit} (tels que $X_{/F}$ soit \inserted{\unsure{ég. loc.|ég.}} cofiltrant donc) $W$-asphériques, donc $X$ est \unsure{tot.|dot.} $W$-asphérique.
% archive sheet number: 70

\scanpage{72}
\handwritten{\diagram{\begin{tikzcd}[ampersand replacement=\&,row sep=small,column sep=small]
 \& {B,C}\unsure{\backslash\!\backslash|\mathbin{/}\mkern-4mu/}\widetilde X\arrow[dl]\arrow[dr,dash] \&\\
 B\backslash\!\backslash\widetilde X\arrow[dr,dash] \& \& C\backslash\!\backslash\widetilde X\arrow[dl,dash]\\
 \& T\arrow[dl,dash] \&\\
 \beta\backslash\!\backslash\widetilde X\arrow[dr,dash] \& \text{cocart} \& \gamma\backslash\!\backslash\widetilde X\arrow[ul]\arrow[dl,dash]\\
 \& \widetilde X\arrow[dl,dash]\arrow[dr] \&\\
 \text{asph.}\ X\arrow[dr] \& \substack{W\\\text{cocart}} \& e\arrow[dl]\\
 \& X' \&
\end{tikzcd}}}
\handwritten{\[
\begin{aligned}
 \widetilde X&\sim\Sigma(T),\quad\text{si }\beta\backslash\!\backslash\widetilde X,\ \gamma\backslash\!\backslash\widetilde X\\
 X'&\sim\Sigma(\widetilde X)\qquad\text{asph. (à vérifier)}
\end{aligned}
\]}
\handwritten{\diagram{\begin{tikzcd}[ampersand replacement=\&,column sep=small]
 \cdot\arrow[r,bend left]\arrow[r,bend right] \& {}\arrow[r,bend left]\arrow[r,bend right] \& {}
\end{tikzcd}}}
\handwritten{\diagram{\begin{tikzcd}[ampersand replacement=\&]
 \alpha\arrow[r,"u'_1"]\arrow[dr,"u''_1"{pos=.3}] \& B\\
 \beta\arrow[ur,"u'_2"{pos=.3}]\arrow[r,"u''_2"'] \& C
\end{tikzcd}}}
\handwritten{\[
\begin{aligned}
 (v u_1,\alpha)&\sim(\unsure{v|o}u_2,\beta)\\
 (v u_1,\alpha)&\sim(v u_2,\beta)
\end{aligned}
\]}
\handwritten{\diagram{\begin{tikzcd}[ampersand replacement=\&,row sep=small]
 \alpha\arrow[dr] \& \\
 \& \unsure{\xi|b}\\
 \beta\arrow[ur] \&
\end{tikzcd}}}
\handwritten{\diagram{\begin{tikzcd}[ampersand replacement=\&,row sep=small]
 \alpha \& b\arrow[dr] \& \unsure{\xi|\zeta}\\
 \unsure{a|\bullet}\arrow[ur,"u_1"]\arrow[rr,bend left,"v"]\arrow[rr,bend right,"w"'] \& \& \unsure{x|n}
\end{tikzcd}}}
\handwritten{\[(v,\unsure{a|\alpha})\sim(w,\beta)\]}
\handwritten{il faut que $v$ se factorise par $u_1$, $w$ par \unsure{$u_2$}.}
\handwritten{\diagram{\begin{tikzcd}[ampersand replacement=\&,row sep=small]
 {}\&\arrow[l,dashed] \unsure{d|\partial}\& \exists \& B\arrow[r,dash] \& \xi\\
 \&\&\& C\arrow[ur] \&
\end{tikzcd}}}
\handwritten{\figure{\begin{tikzpicture}[x=.36cm,y=.36cm]\draw (0,0)--(1.9,11.8);\end{tikzpicture}}}
% letterhead
MINISTERE DE L'EDUCATION NATIONALE\\
ACADEMIE D'AIX-MARSEILLE\\
\rule{2cm}{0.4pt}\\
Inspection Académique\\
de Vaucluse\\
\rule{2cm}{0.4pt}\\
Tél. : 90 82 99 76\\
Référence : Secrétariat n\textsuperscript{o}

AVIGNON, le\\[1ex]
L'Inspecteur d'académie,\\
directeur des services départementaux\\
de l'Education Nationale de Vaucluse,\\[1ex]
à \ M
% archive sheet number: 71

\scanpage{73}\mspage{7}
On voudrait une réciproque, sous des conditions convenables.

\uline{Lemme} $X\in\deleted{\unsure{\mathrm{Ob}}}\mathrm{Cat}$, $F\in X^\wedge$, alors $X_{/F}$ irréd. \unsure{ssi} \inserted{$X_{/F}$ \uline{non vide} et} $\forall$
\diagram{\begin{tikzcd}[ampersand replacement=\&,row sep=small] a\arrow[dr] \& \\ \& F\\ b\arrow[ur] \& \end{tikzcd}}
avec $a,b\in X$, on a $a\times_F b\ne\varnothing$ i.e. \deleted{\unsure{$\exists a$}}
\diagram{\begin{tikzcd}[ampersand replacement=\&,row sep=small] \& a\arrow[dr]\&\\ c\arrow[ur]\arrow[dr,dashed]\&\& F\\ \& b\arrow[ur]\& \end{tikzcd}}
avec $c\in\mathrm{Ob}\,X$, rendant le diagramme comm.

\uline{Corollaire} Pour que $X$ soit \deleted{\unsure{tt|loc.}} \inserted{\deleted{\unsure{loc}} loc.} cofiltrante, \inserted{(i.e. comp. conn. cofiltrantes)} il f. et s. que \deleted{$X$ soit $0$-conn. et que} pour \deleted{$x,y\in$} tout diagramme
% slanted
\margin{\uline{NB} \deleted{\unsure{On doit} supposer ici $X$ connexe,} Il suffit \unsure{prendre|prouver} $\xi=\xi'$.}
\diagram{\begin{tikzcd}[ampersand replacement=\&,row sep=small]
 \& \unsure{\xi|z}\&\\
 a\arrow[ur,"\unsure{\lambda}"]\arrow[dr,"\lambda'"']\&\&b\arrow[ul,"\mu"']\arrow[dl,"\mu'"]\\
 \& \xi' \&\\
 \& x\arrow[uul,dashed,"v'"']\arrow[uur,dashed,"w"] \&
\end{tikzcd}}
dans $X$, on peut compléter en un diagramme commutatif.

Soit $S$ la catégorie ordonnée
\diagram{\begin{tikzcd}[ampersand replacement=\&,row sep=small]
 \&\&\underline{\xi}\arrow[dr]\&\\
 S=\&\underline a\arrow[ur]\arrow[dr]\&\&\underline b,\\
 \&\&\underline{\xi}'\arrow[ur]\&
\end{tikzcd}}
% archive sheet number: 72

\scanpage{74}
% letterhead
MINISTERE DE L'EDUCATION NATIONALE\\
ACADEMIE D'AIX-MARSEILLE\\
\rule{2cm}{0.4pt}\\
Inspection Académique\\
de Vaucluse\\
\rule{2cm}{0.4pt}\\
Tél. : 90 82 99 76\\
Référence : Secrétariat n\textsuperscript{o}

AVIGNON, le\\[1ex]
L'Inspecteur d'académie,\\
directeur des services départementaux\\
de l'Education Nationale de Vaucluse,\\[1ex]
à \ M
\handwritten{\[
\begin{aligned}
 B&=(u_1,\alpha)\sim(u_2,\beta)\\
 C&=(u_2,\alpha)\sim(u_1,\beta)
\end{aligned}
\]}
\handwritten{\diagram{\begin{tikzcd}[ampersand replacement=\&] B\arrow[r]\&\unsure{\xi|\zeta}\end{tikzcd}}}
\handwritten{\diagram{\begin{tikzcd}[ampersand replacement=\&] \alpha\arrow[r,bend left,"u_1"]\arrow[r,bend right,"u_2"']\&b\arrow[r,"v"]\&x\end{tikzcd}}}
\handwritten{\[\xi=(vu_1,\alpha)=(vu_2,\beta)\]}
\handwritten{\diagram{\begin{tikzcd}[ampersand replacement=\&]\deleted{B}C\arrow[r]\&\xi\end{tikzcd}}}
\handwritten{\[\xi=(vu_2,\alpha)=(vu_1,\beta)\]}
\handwritten{\diagram{\begin{tikzcd}[ampersand replacement=\&,row sep=small]
 \alpha\arrow[dr,"\widetilde v"] \&\\ \&\xi\\ \beta\arrow[ur,"\widetilde w"']\&
\end{tikzcd}}}
\handwritten{\diagram{\begin{tikzcd}[ampersand replacement=\&] a\arrow[r,bend left,"v_{\deleted{\unsure{1}}}"]\arrow[r,bend right,"w"']\& \deleted{\unsure{\alpha}}x\end{tikzcd}}}
\handwritten{\[\unsure{\xi_0|\xi}=(v,\alpha)=(w,\beta)\]}
\handwritten{donc $\mathrm{Hom}(b,x)\ne\varnothing$}
\handwritten{\margin{il faut avoir \deleted{\unsure{si l'on}} $v$ se factorise \deleted{que $v$ se factorise} par $u_1$, ou par $u_2$,}}
\handwritten{\figure{\begin{tikzpicture}[x=.35cm,y=.35cm]\draw (0,0)--(3.1,11.8);\end{tikzpicture}}}
% archive sheet number: 73

\scanpage{75}\mspage{8}
$Z$ la catégorie obtenue en ajoutant un objet final (préfabriqué) $e$, d'où
\diagram{\begin{tikzcd}[ampersand replacement=\&]
 \& S\arrow[r,hook,"\text{ouv}"']\arrow[d]\& Z\arrow[d]\&e\arrow[l,hook]\arrow[d,equal]\\
 (D)\&X\arrow[r,hook]\&X'\&e\arrow[l,hook]
\end{tikzcd}}
avec
\[
 \mathrm{Hom}_{X'}(\underset{\substack{\rotatebox{-90}{$\in$}\\ X}}{x},e)\simeq
 \pi_0({}_{x\backslash}S_{/e})\simeq\pi_0({}_{x\backslash}S)
\]
\diagram{\(\displaystyle\simeq\underset{\longrightarrow}{\lim}\left(\vcenter{\hbox{\begin{tikzcd}[ampersand replacement=\&,row sep=small,column sep=small]
 \&\mathrm{Hom}(x,\unsure{\xi|\zeta})\&\\
 \mathrm{Hom}(x,a)\arrow[ur]\arrow[dr]\&\&\mathrm{Hom}(x,b)\arrow[ul]\arrow[dl]\\
 \&\mathrm{Hom}(x,\xi')\&
\end{tikzcd}}}\right)\)}
\[
 \simeq\overbrace{\bigl(\mathrm{Hom}(x,\xi)\amalg\mathrm{Hom}(x,\xi')\bigr)}^{E(x)}/R_{(x)},
\]
$R_{(x)}$ la relation d'équivalence engendrée par les deux relations
\[
 \left\{\begin{aligned}
 \lambda v&\sim\lambda'v&\text{si}\quad v &:x\longrightarrow a\\
 \mu w&\sim\mu'w&\text{si}\quad w &:x\longrightarrow b
 \end{aligned}\right. \ . \tag*{$\unsure{(\alpha)|(a)|(1)|(*)}$}
\]
Le foncteur $S\longrightarrow X$ se remonte canoniquement en $S\longrightarrow X_{/e}$ (grâce au fait que $e$ est objet final de $Z$, cf diagramme $(D)$), soit
% archive sheet number: 74

\scanpage{76}
\handwritten{\[{}_{b\backslash}X\]}
\handwritten{\diagram{\begin{tikzcd}[ampersand replacement=\&]
 {}_{a\backslash}X\arrow[r,"\text{\unsure{cossous|cosous|cossus|cocart.}}"]\&\widetilde X\arrow[r]\&{}_{a\backslash\!\backslash}X\\
 {}_{b\backslash}X\arrow[u,bend left]\arrow[u,bend right]\arrow[ur]\&\&
\end{tikzcd}}}
\handwritten{\diagram{\begin{tikzcd}[ampersand replacement=\&,row sep=small,execute at end picture={\draw (current bounding box.south west)--(current bounding box.north); \draw (current bounding box.south)--(current bounding box.north east);}]
 \&\deleted{b}\arrow[dr,dash]\&\\
 \deleted{a}\arrow[ur,"\deleted{g_1}"]\arrow[dr,"\deleted{g_2}"']\&\&{}\\
 \&\deleted{b}\arrow[ur,dash]\&
\end{tikzcd}}}
\handwritten{\[{}_{\unsure{x|a}\backslash}X\mathbin{\deleted{\unsure{X}}\times}_{\widetilde X}{}_{a\backslash}X\]}
\handwritten{\diagram{\begin{tikzcd}[ampersand replacement=\&]{}_{e\backslash}X\arrow[d]\\X\end{tikzcd}}}
\handwritten{\deleted{\illegible{1}}}
\handwritten{\diagram{\begin{tikzcd}[ampersand replacement=\&,execute at end picture={\draw (current bounding box.west)--(current bounding box.east);}] \deleted{\alpha}\arrow[r]\&\deleted{\beta}\end{tikzcd}}}
\handwritten{\diagram{\begin{tikzcd}[ampersand replacement=\&] \alpha\arrow[r,bend left]\arrow[r,bend right]\&\beta\arrow[r]\&\gamma\end{tikzcd}}}
\handwritten{\diagram{\begin{tikzcd}[ampersand replacement=\&,row sep=small]
 a\arrow[rr,bend left]\arrow[rr,bend right]\arrow[dr]\&\&b\arrow[r]\&c\\
 \&x\arrow[ur]\&\&
\end{tikzcd}}}
\handwritten{\diagram{\begin{tikzcd}[ampersand replacement=\&] \mathrm{Hom}(b,\unsure{b|S|\xi})\arrow[r,shift left=.4ex]\arrow[r,shift right=.4ex]\&\mathrm{Hom}(b,\unsure{a|\bullet|\cdot})\end{tikzcd}}}
\handwritten{\diagram{\begin{tikzcd}[ampersand replacement=\&,row sep=small,execute at end picture={\draw (current bounding box.south west)--(current bounding box.north); \draw (current bounding box.south)--(current bounding box.north east);}]
 \deleted{\alpha}\arrow[r,bend left]\arrow[r,bend right]\&\deleted{B}\arrow[r]\arrow[dr]\&\deleted{C}\\
 \deleted{\unsure{\beta}}\arrow[rr,bend left]\arrow[rr,bend right]\arrow[rr]\&\&\deleted{\unsure{\xi}}
\end{tikzcd}}}
\handwritten{\figure{\begin{tikzpicture}[x=.35cm,y=.35cm]\draw (0,0)--(3.7,11.6); \draw (3.5,1.1)--(6.1,6.7); \draw (4.45,5.7)..controls (4.3,3.6) and (5.3,3.8)..(5.8,3.8)..controls (6.0,3.8) and (6.0,3.6)..(6.25,3.5);\end{tikzpicture}}}
% letterhead
MINISTERE DE L'EDUCATION NATIONALE\\
ACADEMIE D'AIX-MARSEILLE\\
\rule{2cm}{0.4pt}\\
Inspection Académique\\
de Vaucluse\\
\rule{2cm}{0.4pt}\\
Tél. : 90 82 99 76\\
Référence : Secrétariat n\textsuperscript{o}

AVIGNON, le\\[1ex]
L'Inspecteur d'académie,\\
directeur des services départementaux\\
de l'Education Nationale de Vaucluse,\\[1ex]
à \ M
% archive sheet number: 75

\scanpage{77}\mspage{9}
\margin{\unsure{1|$'$}}
\diagram{\begin{tikzcd}[ampersand replacement=\&,row sep=small]
 \&\widetilde{\unsure{\xi|\zeta}}\&\\
 \tilde a\arrow[ur,"\tilde\lambda"]\arrow[dr,"\tilde\lambda'"']\&\&
 \tilde b\arrow[ul,"\tilde\mu"']\arrow[dl,"\tilde\mu'"]\\
 \&\widetilde{\deleted{\unsure{\eta}}\xi}'\&
\end{tikzcd}}
le diagramme correspondant dans $\widetilde X=X_{/e}$.

\uline{Lemme} On a \deleted{$\widetilde X={}_{\tilde\xi\backslash}\widetilde X\cup{}_{\tilde\eta\backslash}\widetilde X$}
\[
\begin{aligned}
 \widetilde X&=\widetilde X_{//\tilde\xi}\cup\widetilde X_{//\unsure{\bar\eta|\tilde\eta|\tilde\xi'}}\\
 \widetilde X_{//\tilde\xi}\cap\widetilde X_{//\tilde\eta}
 &=\widetilde X_{\unsure{\alpha|a}\,//\tilde a}\cup\widetilde X_{//\tilde b}
\end{aligned}
\]
La première relation évidente. Pour la deuxième, soit
\diagram{\begin{tikzcd}[ampersand replacement=\&,row sep=small]
 \&\tilde\xi\&\\ \tilde x\arrow[ur,"\tilde u"]\arrow[dr,dash,"\tilde u'"']\&\&\text{dans}\ \widetilde X\\
 \&\tilde\xi'\&
\end{tikzcd}}
au dessus de
\diagram{\begin{tikzcd}[ampersand replacement=\&,row sep=small]
 \&\xi\&\\x\arrow[ur,"u"]\arrow[dr,"u'"']\&\&\text{dans }X\\\&\unsure{\eta|y}\&
\end{tikzcd}}
% archive sheet number: 76

\scanpage{78}\mspage{\textcircled{6}}
\handwritten{\diagram{\(\displaystyle\overbrace{\vcenter{\hbox{\begin{tikzcd}[ampersand replacement=\&,column sep=small,execute at end picture={\path (leftA.south) -- (rightB.south) node[midway,below] {$\underbrace{\hspace{2.7em}}_{D}$};}]e\arrow[r]\&|[alias=leftA]|a\arrow[r,shift left=.5ex]\arrow[r,shift right=.5ex]\&|[alias=rightB]|b\end{tikzcd}}}}^{Z}\)}}
\handwritten{\diagram{\begin{tikzcd}[ampersand replacement=\&]\bullet\arrow[r,bend left]\arrow[r,bend right]\&\bullet\quad=\end{tikzcd}}}
\handwritten{\diagram{\begin{tikzcd}[ampersand replacement=\&]
 D\arrow[r,hook,"\text{fermé}","\text{\unsure{non ouv}}"' ]\arrow[d]\& Z\arrow[d]\&e\arrow[l,hook,"\text{ouv}"']\arrow[d,equal]\\
 X\arrow[r,hook,"\text{fermé}"]\&X'\&e\arrow[l,hook,"\text{ouv}"']
\end{tikzcd}}}
\handwritten{\[{}_{e\backslash}X_{\deleted{\unsure{/x}}}\]}
\handwritten{\[
\begin{aligned}
 \mathrm{Hom}(e,x)&\simeq\pi_0({}_{e\backslash}D_{/x})\\
 &\simeq\pi_0(D_{/x})\\
 &\simeq\underset{\longrightarrow}{\lim}
 \bigl(\mathrm{Hom}(b,x)\rightrightarrows\mathrm{Hom}(a,x)\bigr)
\end{aligned}
\]}
\handwritten{\diagram{\begin{tikzcd}[ampersand replacement=\&]a\arrow[r,shift left=.4ex]\arrow[r,shift right=.4ex]\&b\end{tikzcd}}}
\handwritten{\deleted{\unsure{$\mathrm{Hom}(b,x)$}}}
\handwritten{un objet de \({}_{e\backslash}X_{\deleted{\unsure{/e}}}\) est déterminé par un objet de $x$, et une flèche $a\longrightarrow x$. L'objet de \({}_{e\backslash}X_{\deleted{\unsure{/e}}}\) défini par $a$, \deleted{\unsure{$\to$}} $a\xrightarrow{\mathrm{id}_a}a$ s'envoie en tous les autres, donc \({}_{e\backslash}X_{\deleted{\unsure{/e}}}\) est $0$-connexe, donc $X\to X'$ est cofinal.}
\handwritten{\figure{\begin{tikzpicture}[x=.36cm,y=.36cm]\draw (0,0)--(3.3,10); \draw (3.65,.1)--(5.3,6.6);\end{tikzpicture}}}
% letterhead
MINISTERE DE L'EDUCATION NATIONALE\\
ACADEMIE D'AIX-MARSEILLE\\
\rule{2cm}{0.4pt}\\
Inspection Académique\\
de Vaucluse\\
\rule{2cm}{0.4pt}\\
Tél. : 90 82 99 76\\
Référence : Secrétariat n\textsuperscript{o}

AVIGNON, le\\[1ex]
L'Inspecteur d'académie,\\
directeur des services départementaux\\
de l'Education Nationale de Vaucluse,\\[1ex]
à \ M
% archive sheet number: 77

\scanpage{79}\mspage{10 //} 
On a donc
\[
\begin{aligned}
 u&\in\mathrm{Hom}(x,\unsure{\xi|\zeta})\subset\mathrm{Hom}(x,\xi)\amalg\mathrm{Hom}(x,\xi')\overset{\rm def}{=}E(x)\\
 u'&\in\mathrm{Hom}(x,\xi')\subset E(\unsure{a|x})\\
 \text{et}\quad u&\sim u'\quad\mathrm{mod}\ R(x).
\end{aligned}
\]
Donc \deleted{il existe} comme $u\ne u'$\inserted{(en tant qu'objets de $E(x)$)}, il existe une chaine d'objets de $E(x)$
\diagram{\begin{tikzcd}[ampersand replacement=\&,row sep=small,column sep=small]
 \&\xi\&\\a\arrow[ur,dash,"\lambda"]\arrow[dr,dash,"\lambda'"']\&\&b\arrow[ul,"\mu"']\arrow[dl,"\mu'"]\\
 \&\xi'\&\\
 x\arrow[uuur,dashed,"u"']\arrow[ur,"u'"]\&\&
\end{tikzcd}}
\[f_0=u,\ f_1,\ f_2,\ \ldots\ f_n=u'\]
$(n\geq1)$, telle que pour $0\leq i\leq n-1$, $f_i$ et $f_{i+1}$ soient reliés par une des deux relations \unsure{$(\alpha)$|$(*)$|$(1)$}\inserted{ou leurs sym.}. En particulier \deleted{\unsure{on aura}} prenant $i=0$, il existe $v:x\to a$ tel que $\underset{\substack{\Vert\\u}}{f_0}=\lambda v$ (et \deleted{\unsure{$\varnothing$}} $f_1=\lambda'v$), ou bien $w:x\to b$ tel que $\underset{\substack{\Vert\\u}}{f_0}=\mu w$, (et $f_1=\mu'w$).

Donc $u$ se factorise soit par $(a,\lambda)$, soit par $(b,\mu)$. Mais cela signifie aussi que $\tilde u:\tilde x\mathrel{\unsure{\leftarrow|-|\to}}\tilde\xi$ se factorise soit par $\tilde a$, soit par $\tilde b$, \deleted{\unsure{qed.}} \deleted{donc que} \deleted{$\tilde\xi$ est} $\tilde\xi$ est dans $\widetilde X_{//\tilde a}\cup\widetilde X_{//\tilde b}$,
qed.
% archive sheet number: 78

\scanpage{80}\mspage{11}
\uline{Cor} $\widetilde X$ est \deleted{$\unsure{0|o}$-}connexe, donc (si les objets $0$-conn. de $X^\wedge$ sont $W$-asph., pour un $W$ donné) il est $W$-asph. (et $X\hookrightarrow X'$ est cocofinal).
\diagram{\begin{tikzcd}[ampersand replacement=\&,row sep=small,column sep=small]
 \&T=\widetilde X_{//(\tilde a,\tilde b)}\arrow[dl]\arrow[dr]\&\\
 \widetilde X_{//\tilde b}\arrow[dr]\&\text{$W$-cocart}\&\widetilde X_{//\tilde a}\arrow[dl]\\
 \&\deleted{\unsure{\widetilde X}}Y\simeq\Sigma^1(T)\arrow[dl]\arrow[dr]\&\\
 \substack{\widetilde X_{//\widetilde{\unsure{\xi|\zeta}}}\\\text{$0$-conn. donc $W$-asph}}\arrow[dr]\&\text{$W$-cocart}\&\substack{\widetilde X_{//\tilde\xi'}\\\text{$0$-conn. car $W$-asph.}}\arrow[dl]\\
 \&\substack{\widetilde X\\ \simeq\Sigma^1_W(Y)}\arrow[dl]\arrow[dr]\&\text{(moyennant $\mathrm{Loc}(6a)$}\\
 \substack{X\\\text{$0$-conn. donc}\\\text{$W$-asph. par hyp}}\arrow[dr,dash]\&\&e\arrow[dl]\\
 \&X'\simeq\Sigma^1_{\deleted{\unsure{W}}}(X)\simeq\Sigma^2_{\deleted{\unsure{W}}}(Y)\simeq\Sigma^3(T)\&
\end{tikzcd}}
Si $T=\varnothing$, alors $X'\overset{W_{\deleted{\unsure{W|\infty}}}}{\sim}\Sigma^2(\unsure{\varepsilon|E})\simeq S^2$ donc si $W$ plus fin que $H^2(\unsure{\overset{''}{-}|-|\cdot},k)$, $k$ \unsure{grpe|grp} comm non \unsure{vide|nul}, donc $X'$ n'est \uline{pas} $W$-asph., absurde. Donc dans ce cas on doit avoir $T\ne\varnothing$
% archive sheet number: 79

\scanpage{81}
% printed form
\noindent\fbox{\begin{minipage}{.94\linewidth}\centering
SEMINAIRE\par\medskip
\textsc{Mediterraneen d'Algebre et Topologie}\par\medskip
M.\,A.\,T\par\medskip
\textit{Deuxième annonce}\par\medskip
\textit{La prochaine session du séminaire} M.\,A.\,T \textit{aura lieu à}\par
\textsc{Toulouse}\par
31 \textsc{mai} et 01 \textsc{juin} 1991\par\medskip
Thème :\par
\textsc{KK-theorie et conjectures de Novikov}\par\medskip
\diagram{\begin{tabular}{@{}l l p{.44\linewidth}@{}}
\multicolumn{3}{l}{Vendredi 31 mai :}\\[.5ex]
10\textsuperscript{h}30 & J.-L. LODAY & \textit{``Les conjectures de Novikov''}.\\
$\left.\begin{array}{c}14^{\mathrm h}15\\ \text{et}\\15^{\mathrm h}30\end{array}\right\}$ & T. FACK & \textit{``Théorie de Kasparov''}.\\
17\textsuperscript{h} & M. WODZICKI & \textit{``Algebraic K-theory and higher index invariants''}.\\[.5ex]
\multicolumn{3}{l}{Samedi 01 juin :}\\[.5ex]
$\left.\begin{array}{c}9^{\mathrm h}\\ \text{et}\\10^{\mathrm h}15\end{array}\right\}$ & G. SKANDALIS & \textit{``Etat de la conjecture de Novikov''}.\\
11\textsuperscript{h}30 & B. TZYGAN & \textit{``Simplicial methods in hermitian K-theory''}.
\end{tabular}}
\textit{Les conférences auront lieu à}\par
UNIVERSITE TOULOUSE III\par
Bâtiment de Mathématiques\par
Salle B - MIG
\end{minipage}}
% archive sheet number: 80

\scanpage{82}\mspage{12}
\deleted{\unsure{$\xi$}}
\diagram{\begin{tikzcd}[ampersand replacement=\&,row sep=small,execute at end picture={\draw (current bounding box.west)--(current bounding box.north)--(current bounding box.center)--(current bounding box.north east)--(current bounding box.south east);}]
 \&\deleted{\tilde a}\\ \deleted{\tilde x}\arrow[ur,dash]\arrow[dr]\&\\ \&\deleted{\tilde b}\end{tikzcd}}
Il faut interpréter $\widetilde X_{//(\tilde a,\tilde b)}$ \unsure{soit|disons} donc
\diagram{\begin{tikzcd}[ampersand replacement=\&,row sep=small,column sep=small]
 \&\&\tilde a\arrow[dl,dashed]\arrow[dr,dashed]\&\\
 \widetilde{\unsure{x|\xi}}\arrow[urr,"v"]\arrow[drr,"w"']\&\widetilde{\unsure{\xi|\zeta}}\&\&\widetilde{\unsure{\xi'|b|c}}\\
 \&\&\tilde b\arrow[ul,dashed]\arrow[ur]\&
\end{tikzcd}}
au dessus de
\diagram{\begin{tikzcd}[ampersand replacement=\&,row sep=small,column sep=small]
 \&\&a\arrow[dl,dashed,"\lambda"']\arrow[dr,dashed,"\lambda'"]\&\\
 x\arrow[urr,"v"]\arrow[drr,"w"']\&\xi\&\&\xi'\\
 \&\&b\arrow[ul,"\mu"]\arrow[ur,"\mu'"']\&
\end{tikzcd}}
Donc $\tilde x$ \deleted{\unsure{$\in E$}} $E(x)/R(x)$ peut être défini indifféremment par
\[
 \underbrace{\lambda v,\mu w}_{\substack{\rotatebox{-90}{$\in$}\\\mathrm{Hom}(x,\xi)}},\quad
 \underbrace{\lambda'v,\mu'w}_{\substack{\rotatebox{-90}{$\in$}\\\mathrm{Hom}(x,\xi')}}
\]
de $E(x)$ \unsure{:|;|\textquotesingle} (\unsure{disons} leurs images dans $E(x)/R(x)$) i.e. ces quatre éléments sont congrus mod $R(x)$. On le sait déjà pour $\lambda v\sim\lambda'v$ et $\mu w\sim\mu'w$, en fait $\tilde x=v^*(\tilde a)=w^{\unsure{\prime\prime|*}}(\tilde b)$, donc il faut exprimer l'égalité
\[\deleted{\unsure{\Leftrightarrow}}\qquad v^*(\tilde a)=w^*(\tilde b)\]
\unsure{ou|donc} encore
\[\boxed{\lambda v\sim\mu w}\]
% archive sheet number: 81

\scanpage{83}
% printed form
\noindent\fbox{\begin{minipage}{.94\linewidth}\centering
SEMINAIRE\par\medskip
\textsc{Mediterraneen d'Algebre et Topologie}\par\medskip
M.\,A.\,T\par\medskip
Session du 31 mai et 01 juin 1991 à \textsc{Toulouse}\par\bigskip
Bulletin d'inscription\par\medskip
\raggedright Nom :\dotfill\par\medskip
Adresse :\dotfill\par\noindent\hbox to \linewidth{\dotfill}\par\medskip
Je désire qu'on me réserve une chambre pour la nuit :\par
\diagram{\(\begin{array}{cc@{\qquad}cc@{\qquad}cc}
 \square &\substack{\text{du 30/05}\\\text{au}\\31/05} &
 \square &\substack{\text{du 31/05}\\\text{au}\\01/06} &
 \square &\substack{\text{du 01/06}\\\text{au}\\02/06}
\end{array}\)}
Bulletin à renvoyer \textbf{le plus rapidement possible} à :\par\medskip
\centering Mme \textsc{Panabiere}\\
\textsc{Universite Toulouse III}\\
\textsc{Topologie et Geometrie}\\
URA 1408\\
118 \textsc{Route de Narbonne}\\
31077 TOULOUSE CEDEX\\
\textsc{Tel.} : 61 55 67 85
\end{minipage}}
% archive sheet number: 82

\scanpage{84}\mspage{13}
\handwritten{Malheureusement, on ne peut en conclure $\lambda v=\mu w$, p. ex. si on a $\lambda'v=\mu'w$, on en conclut $\lambda v\sim\mu w$, mais nullement $\lambda v=\mu w$. Plus précisément, dès qu'on se donne dans $X$ des flèches $v,w$ \deleted{\unsure{avec}}}
\handwritten{\diagram{\begin{tikzcd}[ampersand replacement=\&,row sep=small]
 \&a\\ x'\arrow[ur,"v"]\arrow[dr,"w"']\&\\ \&b
\end{tikzcd}}}
\handwritten{avec $\lambda'v=\mu'\!w$, on en conclut dans $\widetilde X$ un $\tilde x$ avec}
\handwritten{\diagram{\begin{tikzcd}[ampersand replacement=\&,row sep=small]
 \&\tilde a\\ \tilde x\arrow[ur,"\tilde v"]\arrow[dr,"\tilde w"']\&\\ \&\tilde b
\end{tikzcd}}}
\handwritten{, mais on n'aura pas pour \unsure{autant} \deleted{\unsure{on}} $\lambda v=\mu w$ : ainsi, dans le diagramme \deleted{\unsure{de la|de}} \deleted{page} 7, on \inserted{savait déjà qu'on} peut trouver $v,w$ de façon à rendre \uline{un} des carrés pertinents commu (disons $\lambda'v=\mu'w$), \unsure{ce|du|de} qui assure déjà $\widetilde X_{//(\tilde a,\tilde b)}\ne\varnothing$ --}
\handwritten{mais je ne vois pas de moyen d'assurer en $=$ temps que l'autre le soit.}
% letterhead
MINISTERE DE L'EDUCATION NATIONALE\\
ACADEMIE D'AIX-MARSEILLE\\
\rule{2cm}{0.4pt}\\
Inspection Académique\\
de Vaucluse\\
\rule{2cm}{0.4pt}\\
Tél. : 90 82 99 76\\
Référence : Secrétariat n\textsuperscript{o}

AVIGNON, le\\[1ex]
L'Inspecteur d'académie,\\
directeur des services départementaux\\
de l'Education Nationale de Vaucluse,\\[1ex]
à \ M
% archive sheet number: 83

\scanpage{85}
Luc Illusie\\
7 rue Lacaze\\
75014 Paris\\
% archive sheet number: 84

\scanpage{86}
Grothendieck, \unsure{84570|84370} Mormoiron\\
% archive sheet number: 85

